\documentclass[10pt]{article}

\usepackage[utf8]{inputenc}

\usepackage[T1]{fontenc}

\usepackage{amsmath}

\usepackage{amsfonts}

\usepackage{amssymb}

\usepackage[version=4]{mhchem}

\usepackage{stmaryrd}

\usepackage{bbold}



\begin{document}

те КНО-скейлинг не имеет общепризнанного теоретич. объяснения.



М. и. может быть использована для предсказания поведения формфакторов адронов при болыших переданных импульсах и определения структурных функций.



В связи с попытками объяснить в рамках квантовой теории поля скейлинг Бьеркена с начала 1970-х гг. обсуждалась возможность того, что уравнения Дайсона в квантовой теории поля допускают масштабно-инвариантное решение. Для перенормируемой квантовой теории поля этот вопрос оказывается связанным с поведением эффективной константы связности (эффективного заряда) при $-q^{2} \rightarrow \infty$, к-рое определяется видом так наз. функции Гелл-Мана Лоу (см. Ренормализационная группа; дифференциальные уравнения). Для М. и. необходимо, чтобы эта функшия обращалась в нуль при нек-ром значении эффективного заряда. В этом случае при достаточно болыших значениях $-q^{2}$ эффективный заряд совпадает с положением нуля и уравнения ренормализационной группы для вершинных частей обладают масштабно-инвариантными решениями, вообще говоря, с нек-рой аномальной размерностью.



Примеры М. и. с нетривиальными аномальными размерностями имеются в двумерном пространстве-времени (см. Двумерные модели квантовой теории поля). Для перенормируемой квантовой теории поля оказывается, что масштабно-инвариантные решения с необходимостью обладают инвариантностью относительно более общего конформного преобразования, что дает возможность использовать для их нахождения методы конформной квантовой теории поля.



В квантовой хромодинамике асимптотическая свобода приводит к тому, что функция Гелл-Мана - Лоу обращается в нуль при нулевом значении эффективного заряда. В этом случае уравнения ренормализационной группы дают для структурных функций решение, к-рое является функцией не только от отношения $-q_{x}^{2} / M_{x}^{2} c^{2}$, но также слабо (логарифмически) зависит непосредственно от $-q^{2}$. Скейлинг Бъеркена справедлив в квантовой хромодинамике с той точностью, с какой этой дополнительной зависимостью от $-q^{2}$ можно пренебречь. Такое нарушение скейлинга Бьеркена должно наблюдаться в экспериментах по изучению неупругих процессов в достаточно широком диапазоне изменения $-q^{2}$.



Лum.: [1] Боголюбов Н. Н., Ш и рков Д. В., Введение в теорию квантованных полей, 4 изд., М., 1984, гл. 9; [2] С a rruthe rs P., "Phys. Repts», 1971, v. 1С, p. 1; [3] Н икитин Ю. П., Розенталь И. Л., Теория множественных процессов, М., 1976; [4] Дже к и в Р., Приближенная масштабная инвариантность, в кн.: Т рей ман С., Джеки в Р., Гросс Д., Лекции по алгебре токов, пер. с англ., М., 1977, гл. 7. Ю. М. Макеенко. MАСШТАБНОЕ ПРEOБРАЗОВАНИЕ - см. Масштабная инвариантность.



МАТВЕЕВА - ИТСА ФОРМУЛА - формУла для почти периодического по $x$ и по $t$ решения $u(x, t)$ Кортевега - де Фриса уравнения



\[

u_{t}-6 u u_{x}+u_{x x x}=0

\]



Такие решения, зависяшие от $3 l+1$ параметров



\[

E=\left(E_{1}, \ldots, E_{2 l+1}\right), E_{1}<E_{2}<\ldots<E_{2 l+1}, z=\left(z_{1}, \ldots, z_{l}\right),

\]



определяются формулой



\[

u(x, t)=C(E)-\partial^{2} / \partial x^{2} \ln \theta(U x+V t+z, E) .

\]



Здесь $\theta(p, E)=\sum \exp (\operatorname{inp}+i n B(E) n / 2)-l$-мерная $\theta$-функция Римана, $p=\left(p_{1}, \ldots, p_{l}\right)$, суммирование проводится по $l$-мерным целочисленным векторам $n \doteq\left(n_{1}, \ldots, n_{l}\right)$. Элементы $l$-мерных векторов $U=U(E)$ и $V=V(E)$, константа $C(E)$ и элементы $(l \times l)$-симметрич. матрицы $B(E)$ с положительно определенной мнимой частью выражаются через $E$ с



352 МАСШТАБНАЯ помощью гиперэллиптич. интегралов. Аналогичные формулы получены для ряда других нелинейных уравнений.



Лum.: [1] Итс А.Р., Матв е е в В. Б., «Теоретич. и математич. физика», 1975, т. 23, в. 1, с. 51-64; [2] 3 ахаров В. Е., Манак о в С. В., Н о в и ко в С. П., П и т а е в ск и й Л. П., Теория солитоНов, M., 1980.



С. Ю. Доброхотов.



МАТЕМАТИЧЕСКИЙ ВАКУУМ - см. Вакуум.



МАТЕМАТИЧЕСКИЙ МАЯТНИК - см. Маятник.



MATEMАTИЧECKOE OХИДАНИЕ, с р е д н е з наче ние,- одна из числовых характеристик распределения вероятностей случайной величины. Для случайной величины $X$, принимающей последовательность значений $x_{1}, x_{2}, \ldots, x_{k}, \ldots$ с вероятностями, равными соответственно $p_{1}, p_{2}, \ldots, p_{k}, \ldots$, математическое ожидание определяется формулой



\[

\mathrm{E} X=\sum_{k=1}^{\infty} x_{k} p_{k}

\]



при условии, что ряд сходится абсолютно; для случайной величины $X$ с непрерывным распределением, имеющим плотность вероятности $p(x)$,



\[

\mathrm{E} X=\int_{-\infty}^{\infty} x p(x) d x

\]



если интеграл сходится абсолютно. М. о. характеризует расположение значений случайной величины. Полностью эта роль М. о. разъясняется законом болыших чисел. При сложении случайных величин их М. о. складываются, при умножении независимых случайных величин их М.о. перемножаются.



Многие важные характеристики распределений определяются как М. о. нек-рых функций от случайных величин, см., напр., Дисперсия, Момент, Характеристическая функция. В. И. Битюиков. МАТЕМАТИЧЕСКОЙ ФКЗЗКИ КЛАССИЧЕСКИЕ YPABHEHИЯ - уравнения, описывающие математические модели физических явлений. М. ф.к.у. - часть предмета математич. физики. Многие явления физики и механики (гидро- и газодинамики, упругости, электродинамики, оптики, теории переноса, физики плазмы, квантовой физики, теории гравитации и т. п.) описываются краевыми задачами для дифференциальных уравнений. Эти задачи составляют весьма широкий класс М. Ф. к. у.



Для полного описания эволюции физич. процесса помимо уравнений, необходимо, во-первых, задать картину процесса в нек-рый фиксированный момент времени (начальны е условия) и, во-вторых, задать режим на границе той среды, где протекает этот процесс (г р а н и ч ны е усло в ия). Начальные и граничные условия образуют краевые условия, а дифференциальные уравнения вместе с соответствующими краевыми условиями - кра е вы е задач и математич. физики.



Ниже приведены нек-рые примеры уравнений и соответствующих краевых задач.



Уравнение колебаний (волновое уравнение)



\[

\rho \frac{\partial^{2} u}{\partial t^{2}}=\operatorname{div}(p \operatorname{grad} u)-q u+f(x, t)

\]



описывает малые колебания струн, стержней, мембран, акустические и электромагнитные колебания. В уравнении (1) пространственные переменные $x=\left(x_{1}, \ldots, x_{n}\right)$ изменяются в области $G \subset \mathbb{R}^{n}, n=1,2,3$, где развивается рассматриваемый физич. процесс; при этом в соответствии с физич. смыслом входящих величин должно быть $\rho>0, p>0, q>0$.



При $\rho=1, p=a^{2}=$ const и $q=0$ уравнение (1) превращается в волновое уравнение



\[

\frac{\partial^{2} u}{\partial t^{2}}=a^{2} \Delta u+f(x, t)

\]



где $\Delta$ - оператор Лапласа.



Диффузии уравнение



\[

\rho \frac{\partial u}{\partial t}=\operatorname{div}(p \operatorname{grad} u)-q u+f(x, t)

\]





\end{document}